% !TEX root = ../Report.tex
\section{Backgrounds and Related Work}
\label{sec:related_work}

\subsection{Drug--Target Interaction Prediction}

Drug--target interaction (DTI) prediction estimates whether a drug $d$ interacts with a protein target $p$. Each candidate pair is assigned a binary label
\begin{equation}
    y_{dp} \in \{0,1\},
\end{equation}
where $y_{dp}=1$ denotes a known interaction and $y_{dp}=0$ denotes an unknown or sampled negative interaction. A predictive model estimates
\begin{equation}
    \hat{y}_{dp}=F(d,p),
\end{equation}
where $\hat{y}_{dp}$ is the predicted interaction probability. This formulation supports drug repurposing, target identification, virtual screening, and off-target analysis, but it is challenging because verified interactions form only a small fraction of all possible drug--protein pairs.

Two difficulties are especially important. First, supervision is incomplete: unobserved pairs are often sampled as negatives, although some may be unverified positives. Second, the evidence is heterogeneous. Drug structure, protein sequence, pathways, diseases, side effects, and functional annotations provide complementary signals with different scales and semantics. Effective DTI models should therefore integrate molecular and sequence information with relational biomedical context, while also modeling interactions between feature sources.

The experiments use warm-start evaluation, where drugs and proteins in the test pairs also occur in training. The model is evaluated on unseen drug--protein combinations rather than entirely unseen entities. The 1:1 setting contains approximately one sampled negative per positive and provides a balanced discrimination test. The 1:10 setting contains approximately ten sampled negatives per positive, better reflecting DTI sparsity. ROC-AUC and PR-AUC are both reported, with PR-AUC being more informative under the imbalanced 1:10 setting.

\subsection{Biomedical Knowledge Graph Representation}

A biomedical knowledge graph is represented as
\begin{equation}
    \mathcal{G}
    =
    (\mathcal{E},\mathcal{R},\mathcal{T}),
\end{equation}
where $\mathcal{E}$ is the set of entities, $\mathcal{R}$ is the set of relation types, and $\mathcal{T}$ is the set of observed triples. Each triple is written as
\begin{equation}
    (h,r,t) \in \mathcal{T},
\end{equation}
where $h,t\in\mathcal{E}$ are head and tail entities and $r\in\mathcal{R}$ is the relation type. Biomedical entities may include drugs, proteins, genes, diseases, pathways, side effects, biological processes, molecular functions, domains, and protein families. Relations may encode drug--target interactions, drug--disease associations, protein--protein interactions, pathway membership, functional annotations, and domain membership.

Knowledge graph embedding (KGE) maps entities and relations into a continuous vector space. A KGE model defines a scoring function
\begin{equation}
    f(h,r,t)
\end{equation}
that measures the plausibility of a triple. Training encourages observed triples to receive higher scores than implausible triples. For each positive triple $(h,r,t)$, corrupted triples can be generated by replacing the head or tail:
\begin{equation}
    (h,r,t)
    \rightarrow
    (h',r,t),
    \qquad
    (h,r,t)
    \rightarrow
    (h,r,t').
\end{equation}
A common margin-ranking objective separates observed and corrupted triples:
\begin{equation}
    \mathcal{L}_{KGE}
    =
    \sum_{(h,r,t)\in\mathcal{T}^{+}}
    \sum_{(h',r,t')\in\mathcal{T}^{-}}
    \max
    \left(
        0,
        \gamma
        +f(h',r,t')
        -f(h,r,t)
    \right),
    \label{eq:background_kge_loss}
\end{equation}
where $\gamma$ is the margin, $\mathcal{T}^{+}$ contains observed triples, and $\mathcal{T}^{-}$ contains sampled corrupted triples.

Within a DTI framework, KGE compresses graph context into drug and protein representations. For an entity $e$, the KGE stage produces
\begin{equation}
    \mathbf{z}_{e}
    =
    \operatorname{KGE}(e,\mathcal{G}).
\end{equation}
The embeddings $\mathbf{z}_{d}$ and $\mathbf{z}_{p}$ summarize the relational contexts of a candidate drug and protein and can be combined with molecular and sequence descriptors in the downstream classifier. This is useful when direct interaction evidence is sparse but auxiliary biomedical relations remain available.

\subsection{Graph Convolutional Models for Relational Data}

Graph convolutional networks aggregate information from neighboring nodes. For an undirected graph $\mathcal{G}=(\mathcal{V},\mathcal{E},\mathbf{X})$, a standard GCN layer updates node representations as \cite{kipf2017semi}
\begin{equation}
    \mathbf{H}^{k+1}
    =
    \sigma
    \left(
        \widehat{\mathbf{A}}\mathbf{H}^{k}\mathbf{W}^{k}
    \right),
    \qquad
    \mathbf{H}^{0}=\mathbf{X},
    \label{eq:background_gcn_update}
\end{equation}
where $\widehat{\mathbf{A}}=\widetilde{\mathbf{D}}^{-1/2}(\mathbf{A}+\mathbf{I})\widetilde{\mathbf{D}}^{-1/2}$ is the normalized adjacency matrix with self-connections, $\mathbf{W}^{k}$ is a layer-specific transformation, and $\sigma$ is an activation function. Stacking layers lets a node incorporate information from increasingly distant neighborhoods.

Knowledge graphs additionally assign each edge a type and direction. A direct multi-relational extension uses relation-specific transformations:
\begin{equation}
    \mathbf{h}_{v}^{k+1}
    =
    \sigma
    \left(
        \sum_{(u,r)\in\mathcal{N}(v)}
        c_{v,r}^{-1}\mathbf{W}_{r}^{k}\mathbf{h}_{u}^{k}
    \right),
    \label{eq:background_relational_gcn_update}
\end{equation}
where $\mathbf{W}_{r}^{k}$ transforms messages associated with relation $r$, $c_{v,r}$ is a normalization constant, and $\mathcal{N}(v)$ is the relational neighborhood of node $v$. Relation-specific matrices capture edge semantics but increase parameter count as the number of relations grows. Directed-GCN introduces inverse edges and direction-aware sharing \cite{marcheggiani2017encoding}, while R-GCN uses basis or block-diagonal decomposition to control parameters \cite{schlichtkrull2018modeling}. CompGCN instead composes entity and relation embeddings and uses direction-specific transformations, which motivates its use as the stronger KG encoder in Section~\ref{sec:methodology}.

\subsection{Existing DTI Prediction Approaches}

Descriptor-based methods formulate DTI prediction as supervised classification over drug and protein features. Drugs are commonly represented by molecular fingerprints or chemical descriptors, while proteins are represented by sequence-derived descriptors. Classical classifiers and neural networks then learn a decision function from the combined pair representation. These methods are practical and can generalize through structural similarity, but they do not directly exploit heterogeneous biomedical relations.

Network-based methods use the topology of drug, protein, disease, pathway, and other biomedical networks. KG-based methods extend this idea by representing multiple entity and relation types in a unified graph and learning continuous representations that summarize relational compatibility and graph context. Their performance depends on graph coverage, negative sampling, relation modeling, and encoder expressiveness.

Recommendation-based methods interpret drugs and proteins as two interacting entity groups, analogous to users and items. Factorization models learn latent interaction patterns, while neural recommendation models capture nonlinear relationships among heterogeneous feature groups. The KGE-NFM framework combines KG-based and recommendation-based ideas \cite{ye2021unified}: KG embeddings provide relational biomedical representations, descriptor fields provide intrinsic drug and protein information, and NFM models feature interactions among these sources \cite{he2017neural}.
